%%%%%%%%%%%%%%%%%%%%%%%%%%%%%%%%%%%%%%%%%%%%%%%%%%%%%%%%%%%%%%%%%%%%%%%%%%%%%%%%
% conclusion.tex:
%%%%%%%%%%%%%%%%%%%%%%%%%%%%%%%%%%%%%%%%%%%%%%%%%%%%%%%%%%%%%%%%%%%%%%%%%%%%%%%%
\chapter{Conclusion and Discussion}
\label{conclusion_chapter}
%%%%%%%%%%%%%%%%%%%%%%%%%%%%%%%%%%%%%%%%%%%%%%%%%%%%%%%%%%%%%%%%%%%%%%%%%%%%%%%%

The work presented in this dissertation introduces two machine learning algorithms, \aframe and \amplfi, and documents their development, validation, and real-time deployment as an end-to-end gravitational wave search and parameter estimation pipeline.
The central demonstration is that machine learning methods can achieve sensitivity competitive with traditional matched-filter searches for binary black hole mergers while delivering accurate posteriors and skymaps at latencies that are orders of magnitude lower.
Together, \aframe and \amplfi represent the first implementation of an end-to-end machine learning pipeline for gravitational waves operating in a live observing run.

This dissertation also addresses the question of whether such methods are mature enough to be trusted in a production environment alongside the established pipelines that have shaped the field.
Meeting this standard requires more than demonstrating competitive performance in controlled experiments with proof-of-concept implementations.
It requires handling the practical demands of low-latency operation: integration with existing infrastructure, process monitoring and duty cycle requirements, error handling, and real-time human oversight of detected events.
By designing \aframe and \amplfi around these operational requirements and validating them against matched-filtering searches on real observing data, this work provides a template for how machine learning pipelines can be developed and deployed responsibly in gravitational-wave astronomy.

\section{Summary of Contributions}

Chapter~\ref{chap:aframe_methods} introduced \aframe, a convolutional neural network trained to detect compact binary coalescences in two-detector strain data from LIGO.
In contrast to matched-filtering searches, which evaluate a large bank of signal templates, \aframe processes overlapping 1.5-second windows of coincident Hanford and Livingston strain data in real time to produce a scalar detection statistic.
The key technical developments include a robust data augmentation strategy that enables the network to generalize to a wide range of signal morphologies and noise conditions; GPU implementations of common GW data processing tasks, enabling training on large datasets; and an efficient inference-as-a-service deployment capable of analyzing long streams of data on a practical timescale.
We presented the first comparison between machine learning and matched-filtering searches using the same sensitive volume metric reported by the LVK, demonstrating that \aframe achieves competitive sensitivity to binary black hole mergers, while 
also noting that there remain avenues for improving sensitivity to lower-mass systems.
Evaluated on the zero-lag test set from the third observing run, \aframe identified 8 of the 9 GWTC-3 candidates at false alarm rates below 1 per year, the minimum achievable given the duration of timeslides analyzed.
Further, no non-catalog events were reported at false alarm rates below 5 per month, the threshold used for public alerts during O3.

Chapter~\ref{chap:o3_catalog} presented the results of applying \aframe to the full O3 data set.
This analysis recovers 38 events from GWTC-3 with $p_\mathrm{astro} > 0.5$, as well as 2 additional candidates from the IAS catalog and 1 from the OGC catalog, a count which is comparable to that reported by individual matched-filter pipelines operating on the same data.
Of the 15 GWTC-3 candidates not recovered by \aframe at any significance, 5 had secondary masses below $5\,M_\odot$ and fell outside the training prior, while the remaining 10 had network matched-filter SNRs below 12.6.
Parameter estimation was also performed on the IAS and OGC candidates using Bilby with the IMRPhenomXPHM waveform approximant, yielding posteriors broadly consistent with those reported by the respective catalogs.
The O3 search validated the results of Chapter~\ref{chap:aframe_methods} on real data and established \aframe as a viable search pipeline for compact binary mergers, while also highlighting the need to improve sensitivity to lower-mass systems.

Chapter~\ref{chap:amplfi_methods} introduced \amplfi, a likelihood-free inference algorithm for rapid parameter estimation.
\amplfi uses a normalizing flow trained to estimate posterior distributions over source parameters, including chirp mass, mass ratio, and sky location, conditioned on strain data from both LIGO detectors.
We find that \amplfi posteriors are unbiased and the resulting sky maps are similar to those produced by BAYESTAR, the standard low-latency sky localization algorithm, with no significant differences in the areas of 90\% credible regions for well-localized events.
The algorithm produces posterior samples in under one second from the time of a detection trigger, achieving an end-to-end latency of approximately six seconds including detection, significance evaluation, parameter estimation, and sky map generation.
\amplfi was designed for co-deployment with \aframe on a shared GPU, minimizing communication overhead and enabling the full pipeline latency demonstrated here.
The combination of algorithms represents a pipeline capable of participating in the low-latency alert stream alongside traditional matched-filter pipelines. 

Chapter~\ref{chap:amplfi_results} introduced a number of improvements to \amplfi, including the use of multi-modal training data and the inclusion of Virgo data, and  validated it against established methods using both mock data and real data from O3.
Comparisons with BAYESTAR on the mock data stream show that \amplfi produces a similar sky localization area and tighter searched volume constraints.
Additionally, we find that \amplfi posteriors for real events from O3 are broadly consistent with those from Bilby, though less precise, while being orders of magnitude faster.
This includes cases where posterior support extends near or beyond the boundaries of the training prior.
Overall, \amplfi's performance demonstrates its ability to produce data products for low-latency public alerts.

Chapter~\ref{chap:online} described the integration of \aframe and \amplfi into a production pipeline and the deployment of this pipeline during the final segment of LIGO's fourth observing run, O4c, from August 28th to November 25th, 2025.
Of the 24 compact binary candidates for which data from both Hanford and Livingston were available, \aframe contributed to 23 (96\%), and detected 19 (79\%) above the public alert significance threshold.
The single non-detection was due to a power outage that prevented data delivery to the pipeline, rather than a limitation of the algorithm.
For the 9 events where \aframe was the preferred pipeline at the time of the preliminary alert, \amplfi sky maps were distributed to the astronomical community via the General Coordinates Network.
This deployment demonstrates that machine learning methods are sufficiently mature for production use in gravitational wave astronomy and can contribute alongside traditional matched-filter pipelines in an operational setting.

\section{Outlook}

Several extensions of this work are already underway; others represent longer-horizon priorities.

On the detection side, the most significant limitation of the current \aframe implementation is the requirement that both LIGO Hanford and Livingston be operating simultaneously.
Development of multi-detector configurations --- supporting arbitrary detector subsets including three-detector, single-detector, and heterogeneous network modes --- will increase the fraction of observing time for which \aframe can contribute.
An extension of the training mass distribution to lower-mass systems, including binary neutron star and neutron star--black hole mergers, will make the search relevant for the events most likely to have electromagnetic counterparts and yield the most multi-messenger science.
Extending \aframe to BNS systems poses a fundamental signal-duration challenge: a 1.4+1.4\,$M_\odot$ BNS inspiral spends 60--120 seconds in the LIGO sensitive band, far exceeding the 1.5-second input windows on which the current model operates.
One experimental approach under investigation evaluates each input window in the frequency domain against a log-normally spaced grid of chirp mass hypotheses --- approximately 100 templates per window --- using heterodyning to shift each template to baseband, where the slowly evolving signal envelope can be matched efficiently with a compact inner product.
Whether this approach can be made computationally cheap enough for real-time operation remains an open question, but initial experiments are encouraging.

On the parameter estimation side, extending \amplfi to binary neutron star and neutron star--black hole systems is a natural next step given the scientific priority of multi-messenger follow-up.
The longer signal durations and richer parameter spaces of these systems will require dedicated training datasets and may motivate architectural changes to the normalizing flow.
A heterodyning-based approach analogous to that described above for detection is also being explored on the parameter estimation side, compressing long BNS signal representations into a form compatible with \amplfi's existing input window; this work is at an earlier stage than the detection-side development.
Future versions may also incorporate spin-precessing waveform models beyond IMRPhenomD and could support burst signal morphologies, enabling collaboration with unmodeled event search pipelines.

On the deployment side, integrating a distributed task orchestration framework such as Ray or Celery will improve fault tolerance by allowing failed sub-tasks to be retried independently without restarting the entire pipeline.
Several latency reductions have already been made following O4c.
The post-resampling crop was tightened from a conservative 1-second buffer to one based on the filter's transient length (typically less than 0.2 seconds), recovering resampling latency without modifying the underlying filter.
Also in active development is an auxiliary regression head for \aframe to directly output a merger time estimate, which would eliminate the top-hat integration step used to locate the coalescence time.
Longer-horizon improvements include replacing the current symmetric (zero-phase) whitening filter with a causal (minimum-phase) implementation to eliminate the latency introduced by the whitening step, and extending the detector network to include Virgo, which would broaden sky localization coverage for three-detector events.
Infrastructure improvements outside the pipeline will also reduce latency: the next-generation data delivery (NGDD) system, currently in development for future observing runs, will deliver data in 1/16-second chunks rather than 1-second frames, reducing the discrete packet latency, and will deliver each chunk immediately rather than waiting for all auxiliary channels to be ready, reducing the data preparation and transfer latency.

Looking further ahead, the next generation of ground-based detectors --- Cosmic Explorer and the Einstein Telescope --- are expected to detect hundreds of thousands of compact binary mergers per year at far greater distances than current instruments can reach.
At those event rates, machine learning methods will likely be essential infrastructure rather than merely an alternative to matched filtering.
The algorithms, software frameworks, and operational experience developed in this dissertation provide a foundation for meeting that challenge.

%%%%%%%%%%%%%%%%%%%%%%%%%%%%%%%%%%%%%%%%%%%%%%%%%%%%%%%%%%%%%%%%%%%%%%%%%%%%%%%%
